\chapter{有限幾何学}
    \section{有限射影平面}
        \subsection{位数\texorpdfstring{$n$}{n}の有限射影平面には\texorpdfstring{$n^2+n+1$}{n\string^2+n+1}個の点と直線が存在し、各点は\texorpdfstring{$n+1$}{n+1}本の直線に含まれる。}
            \begin{proof}
                \newcommand{\np}{n_\textrm{p}}
                \newcommand{\nl}{n_\textrm{l}}
                \newcommand{\nb}{n_\textrm{b}}
                \quad\par
                念の為、「位数が$n$である」というのは各直線が丁度$n+1$個の点を含むこととして定義されていることを断っておく。
                求めたい点の数を$\np$、直線の数を$\nl$とし、各点が$\nb$本の直線に含まれるとして$\np,\nl,\nb$を求める。
                \par
                まず$\nb$が点に依らず一定であることを示す。
                任意の点$p$を選び、それを含む直線の数を$\nb(p)$とする。
                この$\nb(p)$本の直線によって全ての点がカバーされる。
                なぜならば、もしカバーされない点が存在すれば、有限射影平面の公理1「2つの異なる任意の点が与えられたとき、それらを含むような直線が唯1つだけ存在する。」によりその点と$p$を結ぶ直線が存在して、それは先述の$\nb(p)$本の直線のどれとも一致しない。
                次に公理2「2つの異なる任意の直線の交わりは唯1つの点を含む。」より、$p$は先述の$\nb(p)$本の直線の唯一つの共有点であり、各直線は$p$以外に$n$個の点を含むから
                \[ n \nb(p) + 1 = \np \quad \therefore \nb(p) = (p-1)/n \]
                となり、$\nb(p)$の値は$p$に依存しない。
                この値を$\nb$と表すことにする。
                よって次式が成り立つ。
                \begin{equation}
                    n \nb + 1= \np
                \end{equation}
                \par
                次に、$\nl$本の直線から2本を選ぶ方法の数を2通りに表して方程式を1本立てる。
                1つ目の表し方は明らかに$\binom{\nl}{2}$である。
                公理2より、2本の直線を選ぶと、その共有点として対応する1点が決まる。
                これは一対一対応ではなく、先述のように1点に対してそれを含む直線が$\nb$本存在するので、各点に於いてそれらの直線から2本を選ぶ操作を全ての点に対して行えば、全ての直線から2本の直線を選ぶ方法を網羅できる。
                つまり次式が成り立つ。
                \begin{equation}
                    \binom{\nl}{2} = \np\binom{\nb}{2}
                \end{equation}
                \par
                次に、$\np$個の点から2つを選ぶ方法の数を2通りに表して方程式を1本立てる。
                1つ目の表し方は明らかに$\binom{\np}{2}$である。
                公理1より、2個の点を選ぶと、それらを含む直線が1つ決まる。
                これは一対一対応ではなく、先述のように1直線に対してそれに含まれる点が$n+1$個存在するので、各直線に於いてそれが含む点から2つを選ぶ操作を全ての直線に対して行えば、全ての点から2つの点を選ぶ方法を網羅できる。
                つまり次式が成り立つ。
                \begin{equation}
                    \binom{\np}{2} = \nl\binom{n+1}{2}
                \end{equation}
                式(3)より
                \begin{equation}
                    \nl = \frac{\np(\np-1)}{n(n+1)}
                \end{equation}
                となる。これを式(2)に代入して整理すると
                \[ \frac{\np-1}{n(n+1)}\left(\frac{\np(\np-1)}{n(n+1)} - 1\right) \]
                となる。これに式(1)を適用して整理すると$(\nb-1)(\nb-n-1) = 0$を得る。
                公理3「どの3点も同一直線上にないような4点集合が存在する」より$\nb=1$は不適なので$\nb=n+1$である。
                これと式(1)より$\np = n^2+n+1$である。
                これと式(4)より$\nl = n^2+n+1 = \np$である。
            \end{proof}